\documentclass[11pt]{article}
\usepackage[margin=1in]{geometry}
\usepackage{amsmath,amssymb,amsthm}
\usepackage{booktabs}
\usepackage{graphicx}
\usepackage{natbib}
\usepackage{hyperref}
\newtheorem{theorem}{Theorem}
\newtheorem{lemma}{Lemma}
\newtheorem{proposition}{Proposition}
\newtheorem{corollary}{Corollary}
\newtheorem{conjecture}{Conjecture}
\newtheorem{definition}{Definition}
\newcommand{\sd}{\mathrm{sd}}
\title{The trajectory is not the complexity:\\[2pt]
scan order and the cost of a two-header tape walk}
\author{}
\date{}

\begin{document}
\maketitle

\begin{abstract}
A two-header machine walks on a tape of $n$ cells, cuts at the boundary when a window
statistic fails, and outputs the position where the headers meet. For four such rules the
walk takes $\Theta(n^2)$ steps, and it is natural to read that number as the cost of
computing the output. It is not. An epoch's transition depends on two bits only --- the
statistic at the top window, and whether it fails \emph{anywhere} below it --- and the
second question can be asked from the cheap end, where failures have density one half.
That replaces a scan of the window by an expected constant number of probes, and the
output of all four rules becomes computable at $1.15$ reads per cell, against $48.8$ for
the scan. The quadratic trajectory is real; it is a property of the scan order and not of
the function.

The worst case is different. We prove that the general bound of $\binom{n+2}{3}$ steps is
sharp, with constant $1/6$; that every intermediate exponent between $n$ and $n^3$ is
attained by an explicit rule; and that a random rule is linear. The cost spectrum is
continuous in the worst case and collapses to linear in expectation. We close by asking
what a predictor with sub-linear reads can recover, and state the measurement protocol
that the earlier version of this work lacked.
\end{abstract}

\section{Introduction}

A memoryless walk on a tape is a cheap object to describe and an expensive one to follow.
Two headers sit at the ends of a region, the rule inspects the window between them, and
the action it chooses either moves a header or cuts a cell from an end. The output is the
place where the headers meet. Rules of this shape arise as ``table strategies'' for a
family of automata, and the question that organises the subject is how much of the output
can be had without following the walk.

The answer we were taught to expect is: not much. For four natural rules the walk visits
$\Theta(n^2)$ windows before it stops, and a simulation that follows it costs
$\Theta(n^2)$ reads. The inference from the length of a trajectory to the cost of a
function looks harmless, and it is wrong.

\paragraph{What is actually needed.} Fix the region. If the bit $b_1$ is constant on the
windows an epoch examines, the epoch is a monotone scan of one coordinate: prefixes when
$b_1=0$, suffixes when $b_1=1$. In either case the epoch ends at the \emph{first} window
that fails the second test, and the cut type is determined by three things only: the test
at the top window, whether any window below it fails, and whether the failure is at the
top. The \emph{position} of the failing window is never used. The reference simulation
pays the whole scan to discover that position; a simulator that only asks whether a
failure exists can ask from the bottom, where failures are dense, and pay a constant.

\paragraph{What that changes.} The four rules become linear. At $n=1024$ the scan costs
$48.8$ reads per cell and the reordered evaluation costs $1.15$, a factor of $42$; the
same rule, the same output, the same tapes. The concavity rule whose trajectory we prove
to be quadratic in expectation --- the one that motivated the quadratic picture in the
first place --- is also computable at $1.15$ reads per cell. A trajectory of length
$\Theta(n^2)$ and a computation of cost $\Theta(n)$ coexist, and the gap is entirely in
the order in which the same questions are asked.

\paragraph{What survives of the quadratic.} Two things, and they are worth separating.
First, the trajectory length is a genuine invariant of the rule: for the concavity rule we
prove that the expected number of tests is $\Theta(n^2)$, with an explicit upper bound
$\binom{n+1}{2}$ and a lower bound of $8.1\times 10^{-6}n^2$ that holds without hypothesis.
Second, the worst case over tapes is not linear: on sorted and ramp tapes the reordered
evaluation degenerates to the scan and costs $0.23$ to $0.30\,n^2$. What is false is the
middle term, the identification of the trajectory with the cost of the function.

\paragraph{The spectrum.} Once the two notions are separated, the natural question is
which costs are attainable at all. We show that the general bound $\binom{n+2}{3}$ steps
is sharp --- an explicit rule attains $n^3/6-n^2/2+O(n)$, with a fitted exponent
$3.0098\pm.0024$ --- and that every intermediate exponent can be attained by capping the
construction at $n^{(\beta-1)/2}$ columns. A rule drawn at random, by contrast, is linear:
its field takes an average of $6.8$ steps per region whatever $n$ is. So the worst-case
spectrum is continuous from $n$ to $n^3$ while the expected spectrum collapses to a single
point. Every rule we can construct, including the cubic one, is linear in expectation.

\paragraph{What we do not have.} A ceiling for predictors. An earlier version of this work
reported that a predictor confined to $o(n^2)$ reads recovers a fraction of the output
variance bounded away from one, and read that as evidence for a limit. The budget in those
measurements, $18$ to $50$ reads per cell, is \emph{above} the cost of the reordered
evaluation, which is $1.15$; at those budgets a predictor that uses the reordering is
exact, so the measured fraction describes the scan and not the function. The question is
open and Section~\ref{sec:ceiling} states the protocol that would settle it: budgets from
$0.1\,n$ to $1.2\,n$, several sizes, two-sided intervals, and a correction for the search.

\paragraph{Organisation.} Section~\ref{sec:machine} fixes the machine and the rules.
Section~\ref{sec:order} proves the scan-order theorem and reports the measurements.
Section~\ref{sec:trajectory} separates the trajectory from the cost.
Section~\ref{sec:spectrum} gives the worst-case spectrum.
Section~\ref{sec:ceiling} states what a sub-linear predictor can be asked.
Section~\ref{sec:repro} lists the scripts.

\section{The machine and the rules}\label{sec:machine}

\paragraph{Tape and measure.} $\tau\in\{0,\tfrac12,1\}^n$: the $i$-th cell is the average of
bit $i$ of column one and bit $i$ of column two. Each column is drawn uniformly from the
$\binom nh$ subsets of size $h=\lceil n/2\rceil$, independently, so for even $n$ the cell
law is $P(0)=P(1)=\tfrac14$, $P(\tfrac12)=\tfrac12$, and for odd $n$ it is
$\big(\frac{n-1}{2n}\big)^2$, $\frac{n^2-1}{2n^2}$ and $\big(\frac{n+1}{2n}\big)^2$. Cells
are negatively correlated.

\paragraph{The walk.} A region $[L,R]$ and a window $[\ell,r]\subseteq[L,R]$. At each step
the rule evaluates two bits on the window word and chooses one of eight actions: move
$\ell$ left or right, move $r$ left or right, cut left, cut right, cut both, or halt. A cut
removes the boundary cell of the region; a move that would take a header outside the
region is an \emph{exit} and triggers a cut instead. The run ends in one of three ways: the
headers meet (the output is the meeting position), the region falls below three cells (no
output), or a configuration repeats (no output).

\paragraph{The cut convention.} If the window is the whole region and a move exits, both
ends are cut. If the window is strictly inside the region and a header at an end is asked
to move out, only that end is cut. Ties in the second bit count as false. This convention
is fixed throughout; Section~\ref{sec:repro} records an implementation that got it wrong
and the check that now catches it.

\paragraph{The four rules.} The second bit is the sign of the lag-one autocorrelation,
$\mathrm{corr}(\tau[\ell..r-1],\tau[\ell+1..r])>0$; the first bit is one of: the longest
run of $1$-cells exceeding $m/4$; the number of maximal runs exceeding $m/2$; the presence
of a fixed three-cell pattern; or the window being a palindrome. Write $\rho$ for the
concavity rule, in which the second bit is forced false and the first bit compares the mean
of the first $\lfloor m/2\rfloor$ cells with the mean of all $m$. $\rho$ is the rule whose
trajectory is analysed in Section~\ref{sec:trajectory}.

\paragraph{Cost.} A \emph{read} is one access to one cell. The rule itself reads nothing:
it is a dynamical system driven by the two bits. The machines we cost are the ones that
compute the output, and every access they make --- including the accesses needed to
evaluate a bit --- is charged. Complexity is stated in two regimes, and we keep them apart
throughout: \emph{pathwise}, meaning on every tape of the ensemble, and \emph{in
expectation}, meaning averaged over it.

\section{The scan-order theorem}\label{sec:order}

\subsection{An epoch needs two bits}

Fix the region $[L,R]$ and let $M=R-L+1$.

\begin{lemma}[The epoch is a monotone scan]\label{lem:monotone}
Suppose $b_1=c$ on every window the epoch examines. If $c=0$ then $\ell$ stays at $L$ and
$r$ steps down, so the windows are the prefixes $[L,L+m-1]$ for $m=M,M-1,\dots$. If $c=1$
then $r$ stays at $R$ and $\ell$ steps up, so the windows are the suffixes $[R-m+1,R]$ for
$m=M,M-1,\dots$.
\end{lemma}

\begin{proof}
For $c=0$: $b_2=1$ moves $r$ down by one, which yields the prefix of length $m-1$, and
$b_2=0$ moves $\ell$ to $L-1$, an exit. For $c=1$: $b_2=1$ moves $\ell$ up by one, which
yields the suffix of length $m-1$, and $b_2=0$ moves $r$ to $R+1$, an exit. Induction on
the step index, using that the region is unchanged during the epoch.
\end{proof}

\begin{lemma}[The transition function]\label{lem:transition}
Under the hypothesis of Lemma~\ref{lem:monotone} the cut is a function of
\[
T=b_2\text{ at }m=M,\qquad E=\big[\exists\, m\in[2,M-1]\text{ with }b_2=0\big],
\]
namely: $T=0$ gives a double cut; $T=1$ and $E=1$ gives a single cut on the side the scan
walks towards; $T=1$ and $E=0$ gives a meet.
\end{lemma}

\begin{proof}
By Lemma~\ref{lem:monotone} the epoch walks the prefixes (suffixes) one at a time. $T$
decides whether the first step is already an exit. If it is not, the walk stops at the
largest $m\le M-1$ with $b_2=0$, and the cut type is the same for every such $m$; if there
is none, the walk falls off the bottom and the headers meet. The position of the zero is
never used.
\end{proof}

The content of Lemma~\ref{lem:transition} is that the reference simulation pays $M-m^*+1$
tests to find $m^*$ and then discards it. Only the existence of a failure matters.

\subsection{Certifying that $b_1$ is constant}

Lemma~\ref{lem:transition} applies on a stretch where $b_1$ does not change. For each rule
the lengths at which $b_1$ can deviate are confined to an initial segment $m\le a$, with
$a$ computable after linear preprocessing.

\begin{lemma}[Per-rule certificate]\label{lem:certificate}
Let $c$ be the value of $b_1$ at the top window. Then $b_1=c$ for every $m>a$, where $a$
is: $4\rho$, with $\rho$ the longest run of $1$-cells in the tape, for the longest-run
rule; the longest palindromic prefix at $L$ (Manacher) for the palindrome rule; the
distance to the first occurrence of the pattern at or after $L$, plus two, for the pattern
rule; and the last index at which a $\pm1$ walk $V(j)=2\,\mathrm{ch}(j)-j$ exceeds a
threshold, for the run-count rule.
\end{lemma}

\begin{proof}
Each case is a monotonicity: a long run, a palindrome, an occurrence, or a count that
cannot fall back once it has crossed. For the longest-run rule, $b_1=1$ forces the longest
run of $1$-cells to exceed $m/4$, which is impossible once $m>4\rho$. The other three are
the standard certificates for their statistics.
\end{proof}

Measured, the uncertified zone is small: at $n=1024$ its mean fraction of $M$ is $.081$
for the longest-run rule, $.066$ for run-count, $.062$ for pattern and $.016$ for
palindrome; the epochs that are entirely literal are $1.4$, $0.7$, $1.3$ and $0.1$ per
cent.

\subsection{The algorithm and its cost}

Evaluate $T$; evaluate $E$ by scanning $m=2,3,\dots$ upward, where failures have density
about one half; simulate the tail $m\le a$ literally, which is also where the cycle
detection state lives. Correctness is Lemma~\ref{lem:transition} plus
Lemma~\ref{lem:certificate}; the cost is one probe for $T$, an expected constant number
for $E$, and $O(a)$ for the tail.

\begin{theorem}[Scan order]\label{thm:order}
The output of each of the four rules, and of $\rho$, is computable by an algorithm whose
expected number of reads is linear in $n$ and whose pathwise number is linear on every
tape for which the certificates of Lemma~\ref{lem:certificate} hold throughout.
\end{theorem}

The measurement is in Table~\ref{tab:order}. The reordering is not a heuristic: the
accelerated walk was compared with the reference on every admissible tape for
$n=6,\dots,13$, that is $201{,}643$ tapes at $n=13$, and agrees on the output for all four
rules.

\begin{table}[h]
\begin{center}
\begin{tabular}{lrrrr}
\toprule
& \multicolumn{2}{c}{reference scan} & \multicolumn{2}{c}{reordered} \\
\cmidrule(lr){2-3}\cmidrule(lr){4-5}
$n$ & mean reads & per cell & mean reads & per cell \\
\midrule
$256$  & $3645.9$  & $14.24$ & $298.9$  & $1.168$ \\
$1024$ & $50010.9$ & $48.84$ & $1181.9$ & $1.154$ \\
\bottomrule
\end{tabular}
\caption{Reads to compute the output of the concavity rule $\rho$, mean over tapes and
maximum. At $n=1024$ the reordered evaluation costs a factor $42$ fewer reads; the
pathwise maximum is $3.1$ per cell, against $242$ for the scan.}
\label{tab:order}
\end{center}
\end{table}

\section{Trajectory and cost}\label{sec:trajectory}

The trajectory of $\rho$ is quadratic, and this is a theorem rather than a measurement.
The upper bound is $\binom{n+1}{2}$ tests, since the region shrinks and each epoch tests
at most $M$ windows. The lower bound is $8.1\times10^{-6}n^2$ and holds without hypothesis:
it combines a limit theorem for the first epoch with the frontier inequality, which lets
the epochs be compared with unconditioned suffixes.

So $\rho$ has a trajectory of length $\Theta(n^2)$ and a computation of cost
$\Theta(n)$. The two statements are about different objects and both are true. What the
earlier version of this work did was to identify them.

\section{The worst-case spectrum}\label{sec:spectrum}

Section~\ref{sec:order} showed that a rule can be linear in expectation while its trajectory
is quadratic. This section goes the other way and asks how expensive a rule can be at all,
in the worst case over tapes.

\subsection{The bound is sharp}

\begin{theorem}[General bound]\label{thm:upper2}
Every strategy takes at most $\binom{n+2}{3}$ steps, on every tape.
\end{theorem}

\begin{proof}
Let the segments be $k=0,\dots,K$ with boxes $[L_k,R_k]$ of length $M_k$, and let $P_k$ be the
set of cells visited in segment $k$. Inside a segment the region is fixed, and a repeated
pair $(\ell,r)$ ends the run, so the cells of $P_k$ are distinct and lie in the box:
$|P_k|\le\binom{M_k+1}{2}$. A segment that does not end the run ends with a cut, which
removes at least one cell, so $M_{k+1}\le M_k-1$ and $M_k\le n-k$, with $K\le n$. Hence
\[
\text{steps}=\sum_k|P_k|\le\sum_{k=0}^{n-1}\binom{n-k+1}{2}=\binom{n+2}{3}.
\]
\end{proof}

Two remarks. The bound is largest when the cuts are single, since then $M_k=n-k$ and the
sum is $n^3/6$; double cuts give $n^3/12$. And the bound counts \emph{pairs}: to saturate it
a segment must visit every cell of its box, that is, a Hamiltonian path of the box, and this
must happen for $\Theta(n)$ consecutive boxes.

\subsection{An explicit rule attains $n^3/6$}

\begin{theorem}[Attainment]\label{thm:attain}
There is a strategy whose run takes $\binom{n+2}{3}-n^2+\lfloor n/2\rfloor$ steps on a tape
of its construction, that is $n^3/6-n^2/2+O(n)$.
\end{theorem}

The construction is a column boustrophedon. Start at $c_0=(0,n-1)$, run down column $0$,
advance $\ell$ by one, run up column $1$ to $c_1=(1,n-1)$, advance $\ell$ again, and so on,
so that every corner $c_k=(k,n-1)$ lies on the orbit in order. Every segment ends at the
same last cell $z$, where the field asks for a left cut; the box becomes $[k+1,n-1]$ and the
headers restart exactly at $c_{k+1}$. The field is ``next cell of the orbit'', with a left
cut at $z$ and a halt everywhere else.

Because one field serves every segment, each segment is a suffix of the orbit, and the total
is $\sum_k(|P|-i_k)$ where $i_k$ is the index of the corner $c_k$. Summing the
multiplicities of the cells gives, in closed form,
\[
\text{steps}=\binom{n+2}{3}-n^2+\lfloor n/2\rfloor ,
\]
verified exactly for $n=16,\dots,4096$, with a log--log exponent $3.0098\pm.0024$ over
$n=64,\dots,4096$. The deficit from Theorem~\ref{thm:upper2} is $n^2-\lfloor n/2\rfloor$, and
its two sources are visible in the count: each even column loses two cells at its foot, and
each odd column reaches its corner only at the end, so the cells below the corner carry one
fewer multiplicity. Each contributes $n^2/2$.

\subsection{Every intermediate exponent, and why the natural rules miss them}

Capping the tour at $K=n^{(\beta-1)/2}$ columns gives runs whose fitted exponents are
$2.27\pm.002$, $2.53\pm.003$ and $2.80\pm.003$ for targets $2.25$, $2.5$ and $2.75$. So the
attainable exponents are dense in $(2,3)$, and together with the linear rules of
Section~\ref{sec:order} the worst-case spectrum covers $[1,3]$.

The rules of Section~\ref{sec:machine} do not come near it, and the reason is structural
rather than informational. A rule drawn at random takes an average of $6.8$ steps per
region, independently of $n$, so a random field is linear. What the extremal construction
needs is that the same field drives $\Theta(n)$ segments \emph{and} that the box shrinks by
one cell at a time; the two together cost only $O(n^2)$, which is the deficit of
Theorem~\ref{thm:attain}, and closing it is open. The content channel is not the
obstruction: a window of length $\ge 2\log_3 n$ determines its own position on a random
tape, so $\varphi$ can name a cell once $n$ is large.

\subsection{The spectrum is a worst-case object}

The construction's field is keyed to the tape it runs on, and that tape has probability
about $2^{-n}$ under the ensemble. On any other tape the same field behaves as a random one
and the run is linear. So the continuous spectrum of this section is a statement about
worst case over tapes, and it coexists with the collapse to linearity in expectation
recorded in Section~\ref{sec:order}. Both statements are about the same rules; what differs
is who chooses the tape.

\section{What a sub-linear predictor can recover}\label{sec:ceiling}

\subsection{Three classes}

Sections~\ref{sec:order} and~\ref{sec:spectrum} separate two regimes, and the separation is
worth naming.

\begin{definition}[Classes]\label{def:classes}
Let $g_\rho(n)$ be the least number of cell reads, with $O(\log n)$ bits, that computes
$\delta_\rho$ on every tape, and let $e_\rho(n)$ be the least number that computes it in
expectation. A rule is \emph{linear} if $e_\rho=\Theta(n)$ and $g_\rho=O(n)$; it is
\emph{linear in expectation, superlinear in the worst case} if $e_\rho=\Theta(n)$ and
$g_\rho=\omega(n)$; and it is \emph{superlinear} if $e_\rho=\omega(n)$.
\end{definition}

Every rule measured in this paper is linear or linear in expectation: the canonical rule at
$5.15$ reads per cell, the monotone family at $2$, the four rules of
Section~\ref{sec:machine} at $1.15$, and the constructed rules of
Section~\ref{sec:spectrum} at $1.15$ on random tapes and $n^\beta$ on theirs. The four
rules are of the second kind: on sorted and ramp tapes the reordering degenerates and the
cost is $0.23$ to $0.30\,n^2$.

The third class is empty among everything we can construct, and that is worth stating as a
conjecture rather than leaving as a gap in a definition.

\begin{conjecture}[Expected linearity]\label{conj:linear}
Every strategy has expected cost $e_\rho(n)=O(n)$.
\end{conjecture}

The evidence is that a rule drawn at random takes an average of $6.8$ steps per region,
independently of $n$, so a generic field is linear; the expensive rules of
Section~\ref{sec:spectrum} are expensive only on the tape their field was built for, which
has probability about $2^{-n}$. What the conjecture asserts is that no field is expensive on
a set of tapes of positive measure. It is one-sided: we have an upper bound $1.15n$ on
$e_\rho$ for the four rules and no lower bound, so even $e_\rho=\Theta(n)$ is only half
established.

\subsection{The quantity}

The earlier version of this work reported a ceiling constant of $.56$ from budgets of $18$
to $50$ reads per cell. Those budgets exceed the exact cost of $1.15$ per cell, so a
predictor that uses the reordering is exact there, and the number described the scan rather
than the function; it also conditioned on the tapes the predictor had finished. Both defects
disappear once the budget is placed strictly below the exact cost and the expectation is
taken over the whole output set.

\begin{definition}[Population risk at a sub-linear budget]\label{def:r2}
Fix a rule $\rho$ and $n$, and let $\Omega_n=\{\tau:\delta_\rho(\tau)\neq\bot\}$ with
$\mu_n^\star$ the measure $\mu_n$ conditioned on $\Omega_n$. For a budget $B(n)$ let
$\mathcal P(n,B)$ be the class of adaptive machines with $O(\log n)$ bits of work tape that
read at most $B(n)$ cells and output a value in $\{0,\dots,n-1\}$. Then
\[
R^2(\rho;n,B)\ :=\ 1-\frac{\inf_{P\in\mathcal P(n,B)}\mathbb E_{\mu_n^\star}\big[(P(\tau)-\delta_\rho(\tau))^2\big]}
{\operatorname{Var}_{\mu_n^\star}\big(\delta_\rho(\tau)\big)} .
\]
\end{definition}

\begin{conjecture}[Sub-linear ceiling]\label{conj:ceiling}
Let $\rho$ be linear in expectation and superlinear in the worst case, and let
$B(n)=o(n)$. Then
\[
\limsup_{n\to\infty}\ R^2\big(\rho;\,n,\,B(n)\big)\ \le\ \kappa\ <\ 1
\]
for a constant $\kappa$ that does not depend on $\rho$ or on the choice of $B=o(n)$.
\end{conjecture}

Three comments. The budget is $o(n)$ and not a fixed fraction of $n$: since $e_\rho=\Theta(n)$,
a budget of $cn$ is within a constant factor of the exact cost, a predictor can be exact on a
positive fraction of tapes, and $R^2$ tends to one for reasons that have nothing to do with
the rule --- measured, the curve reaches $.70$ at $B=n/2$ and is contaminated by selection
beyond that point. At $B=o(n)$ the gap to the exact cost diverges, so exactness is
impossible on any tape and the question is quantitative. Second, the conjecture is stated for
the middle class only: for a linear rule a predictor at $B=o(n)$ is short by the same
diverging factor, and the analogous constant is a different number, which we do not claim.
Third, $\kappa$ is a single constant and is the supremum over both arguments,
\[
\kappa=\sup_{\rho}\ \sup_{B=o(n)}\ \limsup_{n}R^2(\rho;n,B),
\]
so a rule that scores lower at a given budget says something about that rule and nothing
about $\kappa$. What the measurements give is a lower bound: if some predictor reaches
$R^2=.56$ at a sub-linear budget, then $\kappa\ge.56$, and the conjecture is that $\kappa<1$.
The bracket $.56\le\kappa<1$ is the honest statement, and the upper end is the part we are
asserting rather than measuring.

\subsection{What a measurement can and cannot establish}

The infimum in Definition~\ref{def:r2} is over all machines and is not computable. A
measurement produces one number: the score of a particular predictor, which is a
\emph{lower} bound on the supremum appearing in the conjecture. Evidence for the conjecture
is that the best predictors we can build fall well short of one; a refutation is a predictor
that does not. Numbers reported in this section must say which side of the inequality they
are on. The earlier version reported a lower bound and called it a ceiling.

\subsection{The exact threshold}

For a per-tape statement the natural object is the decision-tree depth,
\[
D(\delta_\rho)\ :=\ \min\{B:\ \text{some adaptive machine with }B\text{ reads determines }
\delta_\rho\text{ on every }\tau\}.
\]
For $n\le 10$ this is computable exactly, by dynamic programming over the cylinders of the
ensemble: with $E(C,0)=1-\max_v P(\delta=v\mid C)$ and
$E(C,b)=\min_i \mathbb E_v\big[E(C\wedge\{\tau_i=v\},b-1)\big]$, the value $E(\top,B)$ is the
optimal error of an adaptive $B$-read predictor. The computation gives $D(\delta_\rho)=n$ for
the rules tested, that is, exactness requires reading every cell; and it gives the optimal
error at every smaller budget, against which any heuristic can be judged. What it does not
give is the asymptotic law, and in particular whether $D=n$ for all $n$.

\subsection{The protocol}

Budgets $B(n)=n/\log_2 n$, $n^{0.9}$, $n^{0.75}$, $n^{0.5}$; two sizes with $n\ge 256$; a
predictor with $O(\log n)$ bits; the score taken over every output-producing tape with the
unfinished tapes counted as errors; intervals over tapes and over the budget grid; and an
explicit statement, or an explicit refusal, of a correction for the search over rules and
feature sets. The numbers are being measured and are deliberately absent from this draft.

\section{Reproducibility}\label{sec:repro}

Every number comes from a script in the project directory. The scan-order measurements are
in \texttt{probe\_ologn2/} (\texttt{rs\_*.py}, with \texttt{REPORT\_SHORTCUT.md} and its
logs); the trajectory bounds in \texttt{probe\_ologn2/REPORT\_EPOCH\_BOUND.md},
\texttt{REPORT\_UNBALANCED.md} and \texttt{REPORT\_MEET\_BOUND.md}; the spectrum in
\texttt{REPORT\_SUPERQUADRATIC.md}. Tapes are drawn with
\texttt{numpy.random.default\_rng(2026)} unless a paragraph says otherwise.

Two implementation notes. The cut convention of Section~\ref{sec:machine} is the one the
measurements use; an oracle in \texttt{code/bthm0/\_bz\_canon2.py} implemented the single
cut where the double cut is specified, and it has been corrected. The four rules' bits are
evaluated exactly, in integer arithmetic, in every verification reported here.

\bibliographystyle{plainnat}
\begin{thebibliography}{9}
\bibitem[Luo(2026a)]{luo-tww} H.~Luo, \emph{A triangular window walk on random binary
matrices: a conjectured arcsine limit}, manuscript (2026).
\bibitem[Luo(2026b)]{luo-family} H.~Luo, \emph{Memory and randomness control the shape: a
family of tape walks}, manuscript (2026).
\end{thebibliography}

\end{document}
